% ============================================================
%  Chương V – Triển khai và ứng dụng thực tế
% ============================================================
\chapter{Triển khai và ứng dụng thực tế}
\label{chap:deploy}

PKI đi vào hệ thống thực tế bằng một mẫu chung. Một thực thể cần
chứng minh danh tính hoặc quyền phát hành. Thực thể đó dùng chứng
chỉ để gắn danh tính với khóa công khai. Bên nhận kiểm tra chuỗi
chứng chỉ, mục đích sử dụng khóa và trạng thái hiệu lực trước khi
tin vào chữ ký hoặc khóa phiên.

Chương này không lặp lại cấu trúc X.509 đã trình bày ở
Chương~\ref{chap:x509}. Trọng tâm là vai trò vận hành của PKI trong
HTTPS, ký số tài liệu, S/MIME, VPN, ký mã và dịch vụ công điện tử.

\begin{table}[ht]
    \centering
    \caption{So sánh một số ứng dụng PKI}
    \label{tab:pki-applications}
    \scriptsize
    \begin{tabularx}{\textwidth}{@{} >{\raggedright\arraybackslash}p{0.18\textwidth}
        >{\raggedright\arraybackslash}p{0.18\textwidth}
        >{\raggedright\arraybackslash}p{0.18\textwidth}
        X X @{}}
        \toprule
        \textbf{Ứng dụng} & \textbf{Đối tượng xác thực}
            & \textbf{Loại chứng chỉ} & \textbf{Kiểm tra chính}
            & \textbf{Rủi ro vận hành} \\
        \midrule
        HTTPS/TLS
            & Máy chủ dịch vụ
            & TLS server certificate
            & Chuỗi CA, SAN, thời hạn, KU/EKU, thu hồi nếu có
            & Sai hostname, thiếu intermediate CA, khóa máy chủ bị lộ. \\
        Ký số tài liệu
            & Người ký hoặc tổ chức
            & Chứng chỉ ký số cá nhân, tổ chức hoặc dịch vụ ký từ xa
            & Chữ ký tài liệu, chuỗi CA, thời điểm ký, trạng thái thu hồi
            & Không có timestamp tin cậy, khóa ký bị dùng ngoài kiểm soát. \\
        S/MIME
            & Người gửi và người nhận email
            & Chứng chỉ gắn với địa chỉ email
            & Chữ ký email, chứng chỉ người gửi, chứng chỉ người nhận khi mã hóa
            & Khó quản lý chứng chỉ người dùng, mất khóa giải mã. \\
        VPN/IPSec
            & Gateway, client hoặc site ngang hàng
            & Chứng chỉ thiết bị hoặc người dùng
            & CERT, CERTREQ, AUTH trong IKEv2, chuỗi CA nội bộ
            & CA nội bộ yếu, cấp chứng chỉ sai danh tính, thu hồi chậm. \\
        Code signing
            & Nhà phát hành phần mềm
            & Code signing certificate
            & Chữ ký mã, certificate chain, timestamp, trạng thái thu hồi
            & Private key bị lộ, pipeline build bị chiếm quyền. \\
        eGovernment
            & Công dân, tổ chức, cơ quan, dịch vụ ký số
            & Chứng chỉ công cộng hoặc chứng chỉ ký số từ xa
            & Chữ ký văn bản, chứng chỉ người ký, timestamp, OCSP/CRL
            & Liên thông hệ thống, vòng đời chứng chỉ, trải nghiệm người dùng. \\
        \bottomrule
    \end{tabularx}
\end{table}

\begin{figure}[ht]
    \centering
    \begin{tikzpicture}[
        node distance=1.1cm and 1.0cm,
        box/.style={draw, rounded corners, align=center, minimum width=2.5cm,
            minimum height=0.8cm, font=\small},
        arrow/.style={-{Latex[length=2mm]}, thick}
    ]
        \node[box] (root) {Root CA\\trong trust store};
        \node[box, below=of root] (inter) {Intermediate CA};
        \node[box, below=of inter] (leaf) {Chứng chỉ leaf\\của dịch vụ};
        \node[box, right=1.8cm of inter] (client) {Bên kiểm tra};

        \draw[arrow] (root) -- node[right, font=\scriptsize] {ký} (inter);
        \draw[arrow] (inter) -- node[right, font=\scriptsize] {ký} (leaf);
        \draw[arrow] (client) -- node[above, font=\scriptsize] {dựng và kiểm tra chuỗi} (inter);
        \draw[arrow] (client) |- node[pos=0.75, below, font=\scriptsize]
            {kiểm tra tên, mục đích, thời hạn, thu hồi} (leaf);
    \end{tikzpicture}
    \caption{Mẫu kiểm tra PKI trong ứng dụng thực tế}
    \label{fig:pki-validation-flow}
\end{figure}

% ----------------------------------------------------------
\section{HTTPS / TLS}
\label{sec:5.1}

\paragraph{Bài toán thực tế.}
Client cần biết mình đang nói chuyện với đúng máy chủ. Nếu chỉ có
mã hóa kênh truyền, kẻ trung gian vẫn có thể đưa khóa công khai giả.
Trong Web PKI, chứng chỉ máy chủ TLS giải quyết bài toán gắn khóa
công khai với DNS name mà người dùng truy cập.

\paragraph{Cách PKI giải quyết.}
Trong TLS 1.3, máy chủ thường gửi thông điệp \texttt{Certificate}
chứa chứng chỉ leaf và các chứng chỉ trung gian cần thiết. Máy chủ
sau đó dùng khóa bí mật tương ứng để tạo \texttt{CertificateVerify},
qua đó chứng minh quyền kiểm soát khóa trong chứng chỉ
\cite{RFC8446}. Root CA không được gửi như một bằng chứng mới. Nó
phải có sẵn trong trust store của hệ điều hành hoặc trình duyệt.

Certificate chain nối từ leaf đến intermediate CA, rồi đến root CA
được tin cậy. Leaf certificate chứng nhận tên miền dịch vụ.
Intermediate CA giảm rủi ro cho root CA vì root có thể được giữ
offline. Root CA là điểm neo tin cậy mà client dùng để kết thúc quá
trình dựng chuỗi.

\paragraph{Luồng kiểm tra.}
Client kiểm tra chữ ký trên từng mắt xích của chuỗi. Nó từ chối
chứng chỉ nếu thời điểm hiện tại nằm ngoài khoảng hiệu lực. Tên truy
cập phải khớp với DNS name hoặc IP address trong SAN, theo nguyên
tắc đã nêu ở Mục~\ref{sec:4.3}. Client cũng kiểm tra
\texttt{KeyUsage} và \texttt{ExtendedKeyUsage} để bảo đảm chứng chỉ
phù hợp với xác thực máy chủ.

Kiểm tra thu hồi có thể dùng CRL, OCSP hoặc OCSP stapling. OCSP được
RFC~6960 định nghĩa cho truy vấn trạng thái của một chứng chỉ cụ thể
\cite{RFC6960}. OCSP stapling cho phép máy chủ gửi kèm phản hồi OCSP
trong TLS, thay vì để client tự hỏi responder của CA
\cite{RFC6066}. Certificate Transparency bổ sung khả năng giám sát
cấp phát chứng chỉ TLS bằng log công khai. RFC~9162 mô tả CT như một
cơ chế ghi nhận chứng chỉ và precertificate trong log có thể kiểm
toán \cite{RFC9162}.

ACME tự động hóa quá trình yêu cầu, xác minh quyền kiểm soát tên
miền, cấp, gia hạn và thu hồi chứng chỉ. RFC~8555 định nghĩa giao
thức này, thường được biết đến qua hệ sinh thái Let's Encrypt
\cite{RFC8555}. ACME không thay thế kiểm tra PKI ở phía client. Nó
làm giảm lỗi vận hành ở phía cấp và gia hạn chứng chỉ.

\paragraph{Lỗi và rủi ro triển khai.}
Lỗi phổ biến là thiếu intermediate CA trong cấu hình máy chủ. Khi đó
một số client không dựng được chuỗi, dù leaf certificate hợp lệ. Lỗi
khác là chứng chỉ không chứa đúng SAN, hết hạn, dùng sai EKU, hoặc
private key bị sao chép khỏi máy chủ. Trong môi trường nhiều dịch vụ,
tự động hóa gia hạn phải đi kèm giám sát. Một chứng chỉ hết hạn có
thể làm ngừng dịch vụ dù hệ thống ứng dụng vẫn hoạt động.

% ----------------------------------------------------------
\section{Chữ ký số văn bản}
\label{sec:5.2}

\paragraph{Bài toán thực tế.}
Văn bản điện tử cần cơ chế kiểm tra người ký và phát hiện sửa đổi sau
khi ký. Việc chèn hình ảnh chữ ký không cung cấp bảo đảm mật mã. Nó
không chứng minh người ký kiểm soát khóa bí mật, và cũng không ràng
buộc toàn vẹn tài liệu.

\paragraph{Cách PKI giải quyết.}
Luồng ký số bắt đầu bằng việc tính giá trị băm của tài liệu hoặc cấu
trúc dữ liệu cần ký. Phần mềm ký dùng private key để ký giá trị đó.
Chữ ký, chứng chỉ người ký và có thể cả chuỗi chứng chỉ được gắn vào
tài liệu hoặc gói chữ ký. Người nhận dùng public key trong chứng chỉ
để kiểm tra chữ ký, rồi kiểm tra chuỗi CA và trạng thái chứng chỉ.

Chữ ký số, chứng chỉ và timestamp có vai trò khác nhau. Chữ ký số
ràng buộc tài liệu với private key. Chứng chỉ ràng buộc public key
với danh tính người ký hoặc tổ chức. Timestamp ghi nhận rằng chữ ký
đã tồn tại tại một thời điểm do Time Stamping Authority xác nhận.
RFC~3161 định nghĩa Time-Stamp Protocol cho hạ tầng X.509
\cite{RFC3161}.

\paragraph{Luồng kiểm tra.}
Bên nhận kiểm tra chữ ký trên tài liệu trước. Sau đó họ kiểm tra
chứng chỉ người ký, bao gồm chain, thời hạn, mục đích sử dụng khóa và
trạng thái thu hồi. Nếu tài liệu có timestamp đáng tin cậy, bên nhận
có thể đánh giá chữ ký tại thời điểm ký. Điều này giúp kiểm chứng tài
liệu sau khi chứng chỉ người ký đã hết hạn, miễn là chứng chỉ còn hợp
lệ tại thời điểm được đóng dấu.

Với tài liệu PDF, chữ ký số thường được chuẩn hóa theo các profile
riêng của hệ sinh thái tài liệu. Báo cáo này chỉ dùng khái niệm ở mức
PKI: tài liệu được ký, chứng chỉ được đính kèm, và timestamp giúp bảo
vệ giá trị xác minh theo thời gian. Phần chi tiết định dạng PDF không
phải trọng tâm của chương.

\paragraph{Lỗi và rủi ro triển khai.}
Rủi ro chính là bảo vệ private key. Nếu khóa ký nằm trên máy người
dùng mà không có thiết bị bảo vệ hoặc chính sách truy cập phù hợp,
chữ ký có thể bị tạo ngoài ý chí người ký. Timestamp cũng cần nguồn
tin cậy. Nếu tài liệu chỉ có thời gian hệ thống cục bộ, verifier khó
phân biệt chữ ký hợp lệ trước khi hết hạn với chữ ký được tạo muộn
sau đó.

% ----------------------------------------------------------
\section{S/MIME}
\label{sec:5.3}

\paragraph{Bài toán thực tế.}
Email mặc định không cung cấp xác thực mạnh cho nội dung từng thư.
Người nhận cần biết thư có thật do người gửi tạo ra không. Khi cần
bí mật nội dung, người gửi cũng cần mã hóa thư cho đúng người nhận.

\paragraph{Cách PKI giải quyết.}
S/MIME dùng chứng chỉ X.509 gắn với địa chỉ email. RFC~8551 định
nghĩa S/MIME 4.0 cho ký và mã hóa dữ liệu MIME \cite{RFC8551}. Khi
ký email, người gửi dùng private key của mình để tạo chữ ký. Người
nhận dùng chứng chỉ của người gửi để kiểm tra chữ ký và chuỗi CA.

Khi mã hóa email, người gửi cần chứng chỉ của từng người nhận. Nội
dung thư thường được mã hóa bằng khóa phiên. Khóa phiên đó được bảo
vệ bằng public key của người nhận. Chỉ người có private key tương ứng
mới giải mã được nội dung.

\paragraph{Luồng kiểm tra.}
Với thư ký số, client email kiểm tra chữ ký trên nội dung, sau đó
kiểm tra chứng chỉ người gửi. Với thư mã hóa, người nhận dùng private
key của mình để mở khóa phiên, rồi giải mã nội dung. Nếu thư vừa ký
vừa mã hóa, verifier cần xử lý cả hai lớp. Thứ tự xử lý phụ thuộc vào
định dạng thông điệp cụ thể.

S/MIME khác PGP ở mô hình tin cậy. S/MIME dựa vào CA và chứng chỉ
X.509. PGP thường gắn với Web of Trust hoặc cơ chế tự quản lý khóa.
Điều này làm S/MIME phù hợp hơn với tổ chức có PKI tập trung, nhưng
phụ thuộc vào cấp phát và quản trị chứng chỉ người dùng.

\paragraph{Lỗi và rủi ro triển khai.}
Vấn đề thực tế lớn nhất là vòng đời chứng chỉ người dùng. Người dùng
có thể đổi máy, mất private key, hoặc dùng nhiều client email khác
nhau. Nếu private key mã hóa bị mất, thư cũ có thể không đọc lại
được. Nếu tổ chức cấp chứng chỉ cho nhiều địa chỉ hoặc đổi địa chỉ
email, chính sách ánh xạ danh tính cần rõ ràng.

% ----------------------------------------------------------
\section{VPN và IPSec}
\label{sec:5.4}

\paragraph{Bài toán thực tế.}
VPN cần xác thực hai đầu kết nối trước khi thiết lập kênh bảo vệ.
Trong mô hình site-to-site, hai gateway cần nhận ra nhau. Trong mô
hình remote access, gateway cần xác thực client hoặc người dùng. Nếu
chỉ dùng một pre-shared key chung, việc mở rộng và thu hồi quyền truy
cập trở nên khó kiểm soát.

\paragraph{Cách PKI giải quyết.}
IKEv2 thiết lập Security Association cho IPSec. RFC~7296 mô tả việc
các peer có thể dùng payload \texttt{CERT}, \texttt{CERTREQ} và
\texttt{AUTH} trong trao đổi \texttt{IKE\_AUTH} \cite{RFC7296}. Chứng
chỉ cung cấp public key và danh tính. Payload \texttt{AUTH} chứng
minh peer sở hữu private key tương ứng.

Tổ chức thường dùng CA nội bộ cho VPN. CA cấp chứng chỉ cho gateway,
thiết bị hoặc người dùng. Gateway tin root CA nội bộ và kiểm tra
chứng chỉ client theo chính sách truy cập. Hai gateway site-to-site
có thể xác thực lẫn nhau bằng chứng chỉ do cùng một CA hoặc các CA
được liên thông cấp.

\paragraph{Luồng kiểm tra.}
Peer gửi danh tính, chứng chỉ và bằng chứng xác thực. Bên nhận dựng
chuỗi tới CA được tin cậy, kiểm tra thời hạn và trạng thái thu hồi,
rồi kiểm tra rằng danh tính trong chứng chỉ phù hợp với chính sách
VPN. Với mutual authentication, cả hai phía đều thực hiện quy trình
này. Sau khi xác thực thành công, IKEv2 tạo các Child SA để bảo vệ
lưu lượng ESP hoặc AH.

\paragraph{Lỗi và rủi ro triển khai.}
Certificate-based VPN mở rộng tốt hơn pre-shared key vì mỗi thiết bị
hoặc người dùng có thể có chứng chỉ riêng. Khi một client mất quyền,
CA có thể thu hồi chứng chỉ đó mà không đổi bí mật của toàn hệ thống.
Đổi lại, tổ chức phải vận hành PKI nghiêm túc. Lỗi cấu hình CA, bỏ
qua thu hồi, cấp chứng chỉ sai identity hoặc lưu private key gateway
không an toàn đều làm giảm lợi ích của mô hình này.

% ----------------------------------------------------------
\section{Code Signing}
\label{sec:5.5}

\paragraph{Bài toán thực tế.}
Người dùng và hệ điều hành cần biết phần mềm đến từ nhà phát hành nào
và có bị sửa sau khi ký không. Kênh tải xuống không đủ để bảo đảm hai
điều này. Một file có thể bị thay đổi trên mirror, trong quá trình
phân phối, hoặc trong pipeline phát hành.

\paragraph{Cách PKI giải quyết.}
Code signing gắn chữ ký số vào phần mềm, script, driver hoặc gói cài
đặt. Microsoft mô tả chữ ký số trên file như một cách nhận diện nhà
phân phối và bảo đảm nội dung không đổi sau khi chữ ký được tạo
\cite{MicrosoftCryptoTools}. Apple cũng mô tả code signing như cơ chế
xác minh tính toàn vẹn ứng dụng và danh tính nhà phát triển
\cite{AppleCodeSigning}.

Chứng chỉ ký mã khác chứng chỉ TLS ở mục đích sử dụng. Nó không dùng
để chứng minh DNS name của máy chủ. Nó dùng cho publisher hoặc tổ
chức phát hành phần mềm. Sự khác biệt này thường được thể hiện bằng
EKU, chính sách chứng chỉ và quy tắc của nền tảng kiểm tra.

\paragraph{Luồng kiểm tra.}
Khi cài hoặc chạy phần mềm, hệ thống kiểm tra chữ ký trên file. Nếu
chữ ký hợp lệ, hệ thống dựng chuỗi chứng chỉ của publisher tới một
root được tin cậy. Nó cũng có thể kiểm tra timestamp và trạng thái
thu hồi. Timestamp giúp phân biệt phần mềm đã được ký khi chứng chỉ
còn hợp lệ với phần mềm được ký sau khi chứng chỉ hết hạn hoặc bị thu
hồi.

\paragraph{Lỗi và rủi ro triển khai.}
Rủi ro nghiêm trọng nhất là private key ký mã bị lộ. Khi đó kẻ tấn
công có thể ký mã độc dưới danh tính hợp lệ cho tới khi chứng chỉ bị
thu hồi và các nền tảng cập nhật trạng thái. Một rủi ro khác là build
pipeline bị compromise. Trong trường hợp đó, phần mềm độc hại có thể
được ký bằng khóa thật của tổ chức. Vì vậy, code signing cần kiểm soát
khóa, kiểm soát pipeline và quy trình thu hồi rõ ràng.

% ----------------------------------------------------------
\section{eGovernment Việt Nam}
\label{sec:5.6}

\paragraph{Bài toán thực tế.}
Dịch vụ công điện tử cần xác thực công dân, tổ chức, cán bộ xử lý và
văn bản điện tử. Hệ thống cũng cần kiểm tra văn bản đã ký sau khi nó
được gửi qua nhiều cổng dịch vụ hoặc lưu trữ lâu dài. Bài toán chính
không chỉ là tạo chữ ký, mà là quản lý niềm tin ở quy mô quốc gia.

\paragraph{Cách PKI giải quyết.}
Trang thông tin chính thức của Trung tâm Chứng thực điện tử quốc gia
(NEAC) thể hiện các chức năng liên quan đến CA công cộng, ký số tài
liệu, ký số từ xa và kiểm tra văn bản ký số \cite{NEAC}. Ở mức vai
trò hệ thống, CA cấp chứng chỉ cho cá nhân, tổ chức hoặc dịch vụ.
Người dùng dùng chứng chỉ để ký hồ sơ, văn bản hoặc giao dịch. Dịch
vụ công kiểm tra chữ ký và chuỗi chứng chỉ trước khi chấp nhận dữ
liệu.

Ký số từ xa chuyển phần kiểm soát khóa sang dịch vụ được vận hành
chuyên biệt. Người dùng xác thực với dịch vụ ký, yêu cầu ký tài liệu,
rồi nhận chữ ký gắn với chứng chỉ của mình. Mô hình này giảm gánh
nặng cài đặt thiết bị ký ở phía người dùng, nhưng đặt yêu cầu cao hơn
lên xác thực người dùng, nhật ký ký và bảo vệ hệ thống ký.

\paragraph{Luồng kiểm tra.}
Một cổng dịch vụ công hoặc hệ thống văn thư nhận tài liệu đã ký. Nó
kiểm tra chữ ký trên tài liệu, chứng chỉ người ký, chain đến CA được
tin cậy, thời điểm ký và trạng thái thu hồi. Trang NEAC cũng thể hiện
chức năng kiểm tra văn bản PDF/XML, kèm thông tin hiệu lực chứng chỉ,
chữ ký, timestamp, OCSP và CRL \cite{NEAC}. Những trường này phản ánh
các thành phần mà verifier cần khi đánh giá văn bản ký số.

\paragraph{Lỗi và rủi ro triển khai.}
Rủi ro chính là tính liên thông. Các hệ thống khác nhau phải hiểu
đúng định dạng chữ ký, chính sách chứng chỉ, timestamp và trạng thái
thu hồi. Nếu verifier chỉ kiểm tra chữ ký mà bỏ qua chain hoặc thu
hồi, hệ thống có thể chấp nhận chữ ký từ chứng chỉ không còn hợp lệ.
Nếu quy trình ký từ xa xác thực người dùng yếu, chữ ký hợp lệ về mặt
mật mã vẫn có thể không phản ánh đúng ý chí của chủ thể.
